% page 328 du fac-simile (image p-327 du scan)
% bloc 160.0 x 250.6 mm ; marge gauche 8.0 mm
\pgc{-12.700mm}{0.000mm}{
\l{\cel{3}320}
\l{}
\l{\sou{kuboido}. (\sou{anat.}) La mikra osto kubatra, qua esas an la}
\l{\cel{3}parto extera e supra di la tarso. - EF.}
\l{}
\l{\sou{kudo}. I. (\sou{anat.})\cel{2}Angulo salianta quan formacas la parto dopa}
\l{\cel{3}di la artiko di la brakio kun la avanbrakio. - II. Loko}
\l{\cel{3}di la maniko di vesto, qua korespondas a la kubito. -}
\l{\cel{3}III. Angulo salianta (kom ex.: la rivero, la voyo, sinui-}
\l{\cel{3}fas kude). - IV. Parto di utensilo, di instrumento, qua}
\l{\cel{3}formacas angulo salianta. - V. Peco de tubo,formacita}
\l{\cel{3}ek du parti qui facas angulo. - FIS.}
\l{}
\l{\sou{kuglo}. Mikra globo ek metalo, uzata kom projektilo en la}
\l{\cel{3}pafarmi portebla. - D.}
\l{}
\l{kuko. (\sou{koquarto}).Mixuro ek farino, butro, ovi, reduktita a}
\l{\cel{3}pasto, bakita, ed ulakaze sukrizita, aromatizita, gar-}
\l{\cel{3}nisita ye kremo, e c. - DE.}
\l{}
\l{\sou{kukombro}. (bot.) Planto leguma un-yar-vivanta, ek la fami-}
\l{\cel{3}lio \textquotedbl{}kukurbitacei\textquotedbl{}. Frukto di ta planto, oblonga, tre}
\l{\cel{3}aquoza, quan onu manjas o koquita o salade. - L.\cel{1}\sou{cucu}-}
\l{\cel{3}\sou{mis sativus}. - DEFS.}
\l{}
\l{kukulo. (\sou{zool}.) Ucelo klimera, ek la genero \textquotedbl{}pigo\textquotedbl{}. - L.}
\l{\cel{3}\sou{cuculus canorus}. - DEFIRS.}
\l{}
\l{\sou{kukurbito}. (\sou{bot.}) Varietato di gurdego. La frukto di ta}
\l{\cel{3}varietato.- L.\cel{1}\sou{cucurbita}. EFIS.}
\l{}
\l{\sou{kukurbitaceo}. (\sou{bot.}) Familio de planti herbatra, kun sti-}
\l{\cel{3}pito reptera, di qua la gurdego esas la tipo. - EFIS.}
\l{}
\l{\sou{kulo}. l. (\sou{anat.}) La dopa parto (ek la glutei e la supra par-}
\l{\cel{3}to di la kruri) di la korpo homa ed animala.-II. La fun-}
\l{\cel{3}do di ula kozi (botelo, plado, glaso, sako ; en la arti-}
\l{\cel{3}choko : la parto pulpoza qua portas la feno e la foliilFIS}
\l{}
\l{\sou{kulato}. La fundo di la tubo di pafarmo, plu dika kam la ce-}
\l{\cel{3}tero pro ke lu korespondas a la kargajo de pulvero.-FIS.}
\l{}
\l{\sou{kulbutar}. (\sou{netrans}.)I. Bruske falar renverse. - II. Saltar}
\l{\cel{3}rotacante, kapo adavane ed addope. - III. Renversesar}
\l{\cel{3}rotace, tale ke la infrajo divenas la suprajo. - F.}
\l{}
\l{\sou{kulco}. (\sou{zool.}) Mikra moskito. -L.\cel{1}\sou{culex}.}
\l{}
\l{\sou{kulebrino}. Sorto de kanono, longa e mikra-diametra, quan}
\l{\cel{3}onu uzis olim. - EFIS.}
\l{}
\l{\sou{\textquotedbl{}kuli\textquotedbl{}}. En India, Indochinia, Insulo Mauricio : pen-labor-}
\l{\cel{3}isto, portisto. -}
\l{}
\l{kuliero. Implemento di la tablo, di la koqueyo, e c. ek mancho,}
\l{\cel{3}e parto konkava, ovala o cirklatra, per qua onu cherpas, de}
\l{\cel{3}vazo, de plado, e c., nutrivi liquida o mi-liquida, por}
\l{\cel{3}transvarsar li o portar li aden sua boko. - F.}
}
